\section{Introduzione a BlueAgent}\label{introduzione-a-blueagent}

Si vuole presentare dunque il nostro Cyber Range: \emph{BlueAgent}. Un
servizio \emph{HTTP} che implementa delle \emph{API REST} per la
gestione e la verifica delle macchine a partire da file dichiarativi
(\emph{IaC}) fornite dagli organizzatori della CTF. Queste API REST
potranno poi essere sfruttate ed implementate da un qualsiasi
applicativo (denominato CTF Manager) per gestire le challenge, come per
esempio, un plugin per CTFd\footnote{\label{fn:ref_14}\url{https://ctfd.io/}}.

Il Cyber Range implementato da BlueAgent e\textquotesingle{} pensato per
scenari di Patch Management da eseguire su hardware "leggero". Dunque si
vuole offrire un\textquotesingle ambiente basato su container (invece di
utilizzare delle macchine virtuali pesanti) che offra un meccanismo di
verifica completamente nascosto ai giocatori. Tale meccanismo
e\textquotesingle{} basato sulla clonazione sul momento
dell\textquotesingle intera infrastruttura interessata, in una rete
separata e inaccessibile ai giocatori, su cui verranno poi eseguiti i
comandi dell\textquotesingle oracolo definito in fase di dichiarazione.

I membri dei team, potranno connettersi alla rete dei container a loro
assegnati utilizzando una VPN, implementata da WireGuard, tenendo
presente che ogni membro avra associato in WireGuard il suo peer (una
coppia di chiavi pubbliche e private).

E\textquotesingle{} importante puntualizzare che:

\begin{enumerate}
\def\labelenumi{\arabic{enumi}.}
\tightlist
\item
  Gli indirizzi IP dei container, nelle reti dei giocatori, saranno
  statici e noti a priori (questo scenario si applica alla stragrande
  maggioranza delle challenge)
\item
  La verifica dei container, sara implementata attraverso la
  duplicazione della topologia dei container in una rete separata.
  Questa duplicazione avverra sfruttando il meccanismo COW di ZFS.
\item
  BlueAgent non mantiene in memoria lo stato delle challenge (e quindi
  per esempio i punteggi dei vari utenti). Questo software deve essere
  inteso come un\textquotesingle orchestratore di container e non di
  intere challenge. Blueagent si limita a prendere in input la
  definzione di una challenge e di istanziare i container richiesti.
\end{enumerate}

A scopo d\textquotesingle esempio mostriamo un\textquotesingle esempio
per chiarire come creare ed istanziare una challenge con BlueAgent:

\begin{enumerate}
\def\labelenumi{\arabic{enumi}.}
\tightlist
\item
  Il creatore della challenge definisce in modo dichiarativo i
  container.
\item
  Il creatore della challenge, inserisce nel CTF Manager la nuova
  challenge tenendo in memoria tutte le informazioni necessarie a
  BlueAgent (inclusi i file dichiarativi)
\item
  Il giocatore di un determinato team, accede alla challenge e tramite
  il CTF Manager richiede la creazione dei container
\item
  Il CTF Manager invia a BlueAgent i file di configurazione della
  challenge da istanziare
\item
  BlueAgent si occupa di istanziare i container per il team.
\end{enumerate}

BlueAgent e\textquotesingle{} basato su \emph{FastAPI}, per gestire in
modo semplice ed efficiente le richieste HTTP.

L\textquotesingle host che su cui viene eseguito BlueAgent deve mettere
a disposizione dell\textquotesingle architettura le seguenti componenti:

\begin{itemize}
\tightlist
\item
  \textbf{Debian 13 ed il Kernel Linux}: per la gestione delle
  interfacce di rete
\item
  \passthrough{\lstinline!nftables!}: per gestire le regole di inoltro
  nel Kernel
\item
  \textbf{Python 3}: per eseguire BlueAgent.
\item
  \textbf{ZFS}: filesystem che supporta il meccanismo
  \passthrough{\lstinline!COW!}.
\item
  \textbf{OVN ed Open vSwitch}: per gestire la rete di verifica dei
  container.
\item
  \textbf{Distrobuilder}: per generare le immagini dei container in modo
  dichiarativo.
\item
  \textbf{WireGuard}: per fornire una VPN come punto
  d\textquotesingle accesso ai partecipanti
\item
  \textbf{Incus}: per la gestione dei container.
\end{itemize}

Per comprendere passo-passo cosa permette di fare il sistema e come lo
fa sono elencati gli attori, gli use cases ed i requisiti di sistema.
